\documentclass[12pt,a4paper]{article}

\usepackage[T1]{fontenc}
\usepackage[utf8]{inputenc}

\usepackage[spanish]{babel}

\usepackage{amsmath, amssymb, amsthm, mathtools}
\usepackage{esint} 
\usepackage{geometry}
\geometry{margin=2.5cm, top=3cm, bottom=3cm}

\usepackage{xcolor}
\definecolor{azul}{RGB}{0,51,102}
\definecolor{azuloscuro}{RGB}{0,40,80}
\definecolor{grisosuper}{RGB}{51,63,72}
\definecolor{grisclaro}{RGB}{240,242,245}

\usepackage{tikz}
\usetikzlibrary{shapes.geometric, positioning, calc}
\usepackage{pgfplots}
\pgfplotsset{compat=1.18}

\usepackage[most]{tcolorbox}

\usepackage{fancyhdr}
\usepackage{graphicx}
\usepackage{float}
\usepackage{setspace}
\usepackage{titlesec}
\usepackage{booktabs}
\usepackage{enumitem}

\usepackage{listings}
\definecolor{codegreen}{rgb}{0,0.5,0}
\definecolor{codegray}{rgb}{0.5,0.5,0.5}
\definecolor{codepurple}{rgb}{0.58,0,0.82}
\definecolor{backcolour}{rgb}{0.96,0.96,0.94}

\lstdefinestyle{pystyle}{
    backgroundcolor=\color{backcolour},
    commentstyle=\color{codegreen},
    keywordstyle=\color{azul}\bfseries,
    numberstyle=\tiny\color{codegray},
    stringstyle=\color{codepurple},
    basicstyle=\ttfamily\footnotesize,
    breakatwhitespace=false,
    breaklines=true,
    captionpos=b,
    keepspaces=true,
    numbers=left,
    numbersep=6pt,
    showspaces=false,
    showstringspaces=false,
    showtabs=false,
    tabsize=4,
    frame=single,
    rulecolor=\color{gray!40},
    language=Python
}
\lstset{style=pystyle}

\usepackage{hyperref}
\hypersetup{
    colorlinks=true,
    linkcolor=azul,
    urlcolor=azul,
    citecolor=azul
}

\setlength{\headheight}{24.64995pt}
\onehalfspacing

\titleformat{\section}
{\normalfont\fontsize{16}{19}\bfseries\color{grisosuper}}
{\thesection}{1em}{}
[\vspace{0.3em}{\color{azul}\titlerule[2pt]}]

\titleformat{\subsection}
{\normalfont\fontsize{14}{17}\bfseries\color{azul}}
{\thesubsection}{1em}{}

\pagestyle{fancy}
\fancyhf{}
\fancyfoot[C]{\color{grisosuper}\thepage}
\renewcommand{\headrulewidth}{0pt}

\newtcolorbox{cajaconcepto}{
    colback=grisclaro,
    colframe=azul,
    boxrule=2pt,
    arc=3mm,
    left=6pt,
    right=6pt,
    top=6pt,
    bottom=6pt
}

\newtcolorbox{cajaresultado}{
    colback=azul!10,
    colframe=azuloscuro,
    boxrule=2.5pt,
    arc=2mm,
    fonttitle=\bfseries\color{white},
    colbacktitle=azuloscuro,
    title=Resultado
}

\newtcolorbox{cajateorema}{
    colback=azul!5,
    colframe=azul,
    boxrule=2.5pt,
    arc=2mm,
    fonttitle=\bfseries\color{white},
    colbacktitle=azul,
    title=Teorema de Gauss--Ostrogradsky
}

\newtcolorbox{cajadefinicion}[1]{
    colback=grisclaro,
    colframe=grisosuper,
    boxrule=2pt,
    arc=2mm,
    fonttitle=\bfseries\color{white},
    colbacktitle=grisosuper,
    title=#1
}

\begin{document}

\begin{titlepage}
\begin{tikzpicture}[remember picture, overlay]
\fill[azul] (current page.north west) rectangle ($(current page.north east)+(0,-1.5cm)$);
\fill[grisosuper] ($(current page.north west)+(0,-1.5cm)$) --
    ($(current page.north west)+(4cm,-3cm)$) --
    ($(current page.north west)+(4cm,-1.5cm)$) -- cycle;
\fill[azul] (current page.south west) rectangle ($(current page.south east)+(0,1.5cm)$);
\fill[grisosuper] ($(current page.south west)+(0,1.5cm)$) --
    ($(current page.south west)+(4cm,3cm)$) --
    ($(current page.south west)+(4cm,1.5cm)$) -- cycle;
\end{tikzpicture}

\vspace*{3cm}
\centering

{\bfseries\fontsize{20}{24}\selectfont Teorema de Gauss--Ostrogradsky\\
(Teorema de la Divergencia)\par}

\vspace{0.4cm}
{\color{azul}\rule{8cm}{2pt}\par}

\vspace{1cm}

{\large $\displaystyle \iiint_{V} \nabla \cdot \mathbf{F}\, dV \;=\; \oiint_{\partial V} \mathbf{F}\cdot\mathbf{n}\, dS$\par}

\vspace{1.5cm}
{\bfseries Reporte de Investigación\par}

\vspace{2cm}

\begin{tabular}{rl}
\textbf{Nombre del alumno:} & Williams Espinosa López \\
\textbf{Matrícula:} & 251185 \\
\textbf{Cuatrimestre y grupo:} & 4\textordmasculine{} ``B'' \\
\textbf{Docente:} & Sirgei García Ballinas \\
\textbf{Asignatura:} & Cálculo Integral
\end{tabular}

\vfill
\today
\end{titlepage}

\tableofcontents
\newpage

\section{Introducción}

El \textbf{Teorema de Gauss--Ostrogradsky}, también conocido como \textbf{Teorema de la Divergencia}, es uno de los resultados más importantes del cálculo vectorial. Relaciona una integral triple sobre un volumen con una integral de superficie sobre la frontera de ese volumen, y constituye una generalización tridimensional del Teorema Fundamental del Cálculo y del Teorema de Green en el plano.

Históricamente fue formulado de manera independiente por el matemático alemán \textbf{Carl Friedrich Gauss} (1777--1855) y por el matemático ruso \textbf{Mijaíl Vasílievich Ostrogradsky} (1801--1862), este último publicando la demostración general en 1826.

\begin{cajaconcepto}
\textbf{Objetivo:} Estudiar el Teorema de Gauss--Ostrogradsky desde su formulación teórica y verificarlo analítica y numéricamente mediante tres ejemplos prácticos implementados en Python, comprobando la igualdad entre la integral de volumen de la divergencia y el flujo del campo a través de la superficie frontera.
\end{cajaconcepto}

El teorema tiene aplicaciones fundamentales en electromagnetismo (ley de Gauss), mecánica de fluidos (ecuación de continuidad), transferencia de calor, gravitación y ecuaciones diferenciales parciales.

\section{Fundamentos Teóricos}

\subsection{Conceptos previos}

\begin{cajadefinicion}{Definición 1: Campo vectorial}
Un \textbf{campo vectorial} en $\mathbb{R}^3$ es una función $\mathbf{F}:\mathbb{R}^3 \to \mathbb{R}^3$ que asigna a cada punto $(x,y,z)$ un vector:
\[
\mathbf{F}(x,y,z) = \big(P(x,y,z),\; Q(x,y,z),\; R(x,y,z)\big).
\]
\end{cajadefinicion}

\begin{cajadefinicion}{Definición 2: Divergencia}
La \textbf{divergencia} de un campo vectorial $\mathbf{F}=(P,Q,R)$ de clase $C^1$ es el campo escalar:
\[
\nabla \cdot \mathbf{F} \;=\;
\frac{\partial P}{\partial x} +
\frac{\partial Q}{\partial y} +
\frac{\partial R}{\partial z}
\]
Mide la razón neta a la que el campo emana de un punto: si $\nabla\cdot\mathbf{F} > 0$ el punto actúa como \emph{fuente}, si $\nabla\cdot\mathbf{F} < 0$ actúa como \emph{sumidero}.
\end{cajadefinicion}

\begin{cajadefinicion}{Definición 3: Flujo de un campo vectorial}
El \textbf{flujo} de un campo vectorial $\mathbf{F}$ a través de una superficie orientada $S$ con vector normal unitario $\mathbf{n}$ es:
\[
\Phi \;=\; \iint_{S} \mathbf{F}\cdot\mathbf{n}\, dS.
\]
Representa la cantidad neta del campo que atraviesa $S$ en la dirección de $\mathbf{n}$.
\end{cajadefinicion}

\subsection{Enunciado del teorema}

\begin{cajateorema}
Sea $V\subset \mathbb{R}^3$ un sólido acotado cuya frontera $\partial V = S$ es una superficie cerrada, regular a trozos, orientada hacia el exterior mediante el vector normal unitario $\mathbf{n}$. Sea $\mathbf{F}$ un campo vectorial de clase $C^1$ definido sobre un abierto que contiene a $V$. Entonces:
\[
\iiint_{V} (\nabla \cdot \mathbf{F}) \, dV \;=\;
\oiint_{S} \mathbf{F}\cdot \mathbf{n}\, dS
\]
\end{cajateorema}

\subsection{Interpretación física}

El teorema afirma que \emph{la cantidad total de ``materia'' que un campo vectorial genera en el interior de un volumen es igual a la cantidad neta que atraviesa su frontera}. En mecánica de fluidos, si $\mathbf{F}$ es el campo de velocidades, la divergencia mide las fuentes/sumideros puntuales, mientras que el flujo mide cuánto fluido sale efectivamente por la superficie. El teorema garantiza que ambos totales deben coincidir.

\subsection{Condiciones de aplicabilidad}

Para aplicar correctamente el teorema se deben cumplir:
\begin{enumerate}[leftmargin=1.5cm]
    \item $V$ es un dominio sólido acotado y simplemente conexo.
    \item La frontera $\partial V$ es cerrada y regular a trozos.
    \item El vector normal $\mathbf{n}$ apunta hacia el exterior del volumen.
    \item $\mathbf{F}$ tiene derivadas parciales continuas en un abierto que contiene $V\cup\partial V$.
\end{enumerate}

\subsection{Fórmulas en distintos sistemas de coordenadas}

\textbf{Coordenadas cartesianas:} $dV = dx\, dy\, dz$
\[
\nabla\cdot\mathbf{F} = \frac{\partial P}{\partial x} + \frac{\partial Q}{\partial y} + \frac{\partial R}{\partial z}.
\]

\textbf{Coordenadas cilíndricas:} $dV = r\, dr\, d\theta\, dz$
\[
\nabla\cdot\mathbf{F} = \frac{1}{r}\frac{\partial(r F_r)}{\partial r} + \frac{1}{r}\frac{\partial F_\theta}{\partial \theta} + \frac{\partial F_z}{\partial z}.
\]

\textbf{Coordenadas esféricas:} $dV = \rho^2 \sin\varphi\, d\rho\, d\varphi\, d\theta$
\[
\nabla\cdot\mathbf{F} = \frac{1}{\rho^2}\frac{\partial(\rho^2 F_\rho)}{\partial \rho} + \frac{1}{\rho\sin\varphi}\frac{\partial(\sin\varphi\, F_\varphi)}{\partial \varphi} + \frac{1}{\rho\sin\varphi}\frac{\partial F_\theta}{\partial \theta}.
\]

\section{Metodología}

Para cada uno de los tres ejemplos se procedió de la siguiente manera:
\begin{enumerate}[leftmargin=1.5cm]
    \item Se definió el campo vectorial $\mathbf{F}$ y la región $V$.
    \item Se calculó analíticamente $\nabla\cdot\mathbf{F}$ y se evaluó la integral triple $\iiint_V \nabla\cdot\mathbf{F}\, dV$.
    \item Se calculó analíticamente el flujo $\oiint_{\partial V} \mathbf{F}\cdot\mathbf{n}\, dS$.
    \item Se implementó en Python la verificación numérica de ambos lados usando \texttt{scipy.integrate}.
    \item Se generaron gráficas 3D con \texttt{matplotlib} para visualizar el campo, la región y la comparación numérica.
\end{enumerate}

\section{Ejemplo 1: Campo radial sobre una esfera}

\subsection{Planteamiento}
Sea $\mathbf{F}(x,y,z)=(x,\,y,\,z)$ y $V$ la bola sólida $x^2+y^2+z^2\le R^2$ con $R=1$. Verificar el teorema de Gauss--Ostrogradsky.

\subsection{Lado izquierdo: integral de la divergencia}

\[
\nabla\cdot\mathbf{F} = \frac{\partial x}{\partial x} + \frac{\partial y}{\partial y} + \frac{\partial z}{\partial z} = 3.
\]

Por lo tanto:
\[
\iiint_V \nabla\cdot\mathbf{F}\, dV
= 3 \iiint_V dV
= 3 \cdot \frac{4}{3}\pi R^3 = 4\pi R^3.
\]

\subsection{Lado derecho: flujo a través de la esfera}

En la esfera el vector normal unitario exterior es $\mathbf{n} = (x,y,z)/R$. Entonces:
\[
\mathbf{F}\cdot\mathbf{n} = \frac{x^2+y^2+z^2}{R} = \frac{R^2}{R} = R.
\]

Por lo tanto:
\[
\oiint_{\partial V} \mathbf{F}\cdot\mathbf{n}\, dS
= R \cdot \text{Área}(\partial V)
= R \cdot 4\pi R^2 = 4\pi R^3.
\]

\begin{cajaresultado}
Ambos lados coinciden:
\[
\iiint_V \nabla\cdot\mathbf{F}\, dV = \oiint_{\partial V} \mathbf{F}\cdot\mathbf{n}\, dS = 4\pi R^3.
\]
Para $R=1$: $4\pi \approx 12.5664$. \hspace{0.5cm}$\checkmark$\; Teorema verificado.
\end{cajaresultado}

\begin{figure}[H]
    \centering
    \includegraphics[width=0.95\textwidth]{ejemplo1_esfera.png}
    \caption{Izquierda: campo radial $\mathbf{F}=(x,y,z)$ sobre la esfera unitaria. Los vectores son perpendiculares a la superficie, apuntando hacia afuera. Derecha: comparación numérica de ambos lados del teorema.}
\end{figure}

\section{Ejemplo 2: Campo cuadrático sobre un cubo}

\subsection{Planteamiento}
Sea $\mathbf{F}(x,y,z) = (x^2,\,y^2,\,z^2)$ y $V$ el cubo unitario $[0,a]^3$ con $a=1$. Verificar el teorema.

\subsection{Lado izquierdo: integral de la divergencia}

\[
\nabla\cdot\mathbf{F} = 2x + 2y + 2z.
\]

\begin{align*}
\iiint_V (2x+2y+2z)\, dV
&= 2\int_0^a\!\!\int_0^a\!\!\int_0^a x\, dx\,dy\,dz + \cdots\\
&= 2\cdot \frac{a^2}{2}\cdot a^2 \;\cdot\; 3 \;=\; 3a^4.
\end{align*}

Para $a=1$: $3a^4 = 3$.

\subsection{Lado derecho: flujo por las 6 caras del cubo}

\begin{table}[H]
\centering
\begin{tabular}{lccc}
\toprule
\textbf{Cara} & \textbf{Normal} $\mathbf{n}$ & $\mathbf{F}\cdot\mathbf{n}$ & \textbf{Flujo} \\
\midrule
$x=a$  & $(1,0,0)$  & $x^2 = a^2$   & $a^2 \cdot a^2 = a^4$ \\
$x=0$  & $(-1,0,0)$ & $-x^2 = 0$    & $0$ \\
$y=a$  & $(0,1,0)$  & $y^2 = a^2$   & $a^4$ \\
$y=0$  & $(0,-1,0)$ & $-y^2 = 0$    & $0$ \\
$z=a$  & $(0,0,1)$  & $z^2 = a^2$   & $a^4$ \\
$z=0$  & $(0,0,-1)$ & $-z^2 = 0$    & $0$ \\
\midrule
\multicolumn{3}{r}{\textbf{Flujo total:}} & $3a^4$ \\
\bottomrule
\end{tabular}
\caption{Flujo a través de cada cara del cubo.}
\end{table}

\begin{cajaresultado}
Ambos lados dan $3a^4 = 3$ para $a=1$. \hspace{0.5cm}$\checkmark$\; Teorema verificado.
\end{cajaresultado}

\begin{figure}[H]
    \centering
    \includegraphics[width=\textwidth]{ejemplo2_cubo.png}
    \caption{Izquierda: cubo unitario con el campo $\mathbf{F}=(x^2,y^2,z^2)$. Derecha: flujo por cada cara; solo las caras en $+x$, $+y$, $+z$ contribuyen, y su suma iguala a la integral de la divergencia.}
\end{figure}

\section{Ejemplo 3: Campo sobre un cilindro}

\subsection{Planteamiento}
Sea $\mathbf{F}(x,y,z) = (xz,\,yz,\,z^2)$ y $V$ el cilindro $x^2+y^2 \le R^2$, $0\le z\le h$, con $R=1$, $h=2$.

\subsection{Lado izquierdo: integral de la divergencia}

\[
\nabla\cdot\mathbf{F} = z + z + 2z = 4z.
\]

En coordenadas cilíndricas, con $dV = r\, dr\, d\theta\, dz$:
\begin{align*}
\iiint_V 4z\, dV
&= \int_0^{2\pi}\!\!\int_0^R\!\!\int_0^h 4z \cdot r\, dz\, dr\, d\theta \\
&= 2\pi \cdot \frac{R^2}{2} \cdot 2h^2
\;=\; 2\pi R^2 h^2.
\end{align*}

Para $R=1$, $h=2$: $2\pi(1)(4) = 8\pi \approx 25.1327$.

\subsection{Lado derecho: flujo por las tres superficies}

\textbf{(a) Tapa superior} $z=h$, $\mathbf{n}=(0,0,1)$:
\[
\mathbf{F}\cdot\mathbf{n} = z^2 = h^2
\;\;\Rightarrow\;\;
\Phi_{\text{sup}} = h^2 \cdot \pi R^2.
\]

\textbf{(b) Tapa inferior} $z=0$, $\mathbf{n}=(0,0,-1)$:
\[
\mathbf{F}\cdot\mathbf{n} = -z^2 = 0
\;\;\Rightarrow\;\;
\Phi_{\text{inf}} = 0.
\]

\textbf{(c) Superficie lateral}, $\mathbf{n}=(x/R,\,y/R,\,0)$:
\[
\mathbf{F}\cdot\mathbf{n} = \frac{xz\cdot x + yz\cdot y}{R}
= \frac{z(x^2+y^2)}{R} = \frac{z\,R^2}{R} = zR.
\]
Con $dS = R\, dz\, d\theta$:
\[
\Phi_{\text{lat}} = \int_0^{2\pi}\!\!\int_0^h zR \cdot R\, dz\, d\theta
= 2\pi R^2 \cdot \frac{h^2}{2} = \pi R^2 h^2.
\]

\textbf{Flujo total:}
\[
\Phi = \pi R^2 h^2 + \pi R^2 h^2 + 0 = 2\pi R^2 h^2.
\]

\begin{cajaresultado}
Ambos lados dan $2\pi R^2 h^2 = 8\pi \approx 25.1327$ para $R=1$, $h=2$. \hspace{0.5cm}$\checkmark$\; Teorema verificado.
\end{cajaresultado}

\begin{figure}[H]
    \centering
    \includegraphics[width=\textwidth]{ejemplo3_cilindro.png}
    \caption{Izquierda: cilindro con el campo $\mathbf{F}=(xz,yz,z^2)$. Derecha: contribuciones al flujo por la tapa inferior, superficie lateral y tapa superior, cuya suma iguala a la integral triple de la divergencia.}
\end{figure}

\section{Resumen comparativo}

\begin{table}[H]
\centering
\begin{tabular}{lccc}
\toprule
\textbf{Región} & \textbf{Campo $\mathbf{F}$} &
$\iiint_V \nabla\cdot\mathbf{F}\, dV$ &
$\oiint_{\partial V} \mathbf{F}\cdot\mathbf{n}\, dS$ \\
\midrule
Esfera ($R=1$)           & $(x,y,z)$       & $12.5664$ & $12.5664$ \\
Cubo ($a=1$)             & $(x^2,y^2,z^2)$ & $3.0000$  & $3.0000$  \\
Cilindro ($R=1, h=2$)    & $(xz,yz,z^2)$   & $25.1327$ & $25.1327$ \\
\bottomrule
\end{tabular}
\caption{En los tres casos, la integral de volumen y el flujo de superficie coinciden numéricamente, confirmando el teorema.}
\end{table}

\begin{figure}[H]
    \centering
    \includegraphics[width=0.85\textwidth]{resumen_comparativo.png}
    \caption{Comparación global de los tres ejemplos. Cada par de barras (volumen vs superficie) tiene la misma altura, lo que verifica visualmente el Teorema de Gauss--Ostrogradsky.}
\end{figure}

\section{Conclusiones}

\begin{cajaconcepto}
El Teorema de Gauss--Ostrogradsky fue verificado analítica y numéricamente en tres geometrías distintas (esfera, cubo y cilindro), con tres campos vectoriales diferentes. En los tres casos la diferencia entre los dos lados del teorema fue prácticamente nula, del orden del error de máquina.
\end{cajaconcepto}

\begin{itemize}[leftmargin=1.5cm]
    \item El teorema permite \emph{reducir} una integral de superficie (a menudo tediosa, sobre una frontera posiblemente compuesta por varias piezas) a una integral de volumen (normalmente más simple). En el ejemplo del cubo, por ejemplo, es más fácil integrar $2x+2y+2z$ sobre el cubo que evaluar el flujo por cada una de las seis caras.
    \item Las condiciones de aplicabilidad (región cerrada, frontera regular, campo $C^1$, normal saliente) son esenciales: violar alguna de ellas puede conducir a resultados incorrectos.
    \item La implementación numérica con \texttt{scipy.integrate} confirma que las cuadraturas adaptativas de Python son herramientas confiables para verificar teoremas del cálculo integral con precisión prácticamente exacta.
    \item Este teorema es la base teórica de resultados físicos fundamentales como la \emph{ley de Gauss} del electromagnetismo y la \emph{ecuación de continuidad} en mecánica de fluidos.
\end{itemize}


\end{document}
